%!TEX root = ../main.tex
% Declaración sobre el uso de inteligencia artificial.

Durante la elaboración del presente trabajo final, titulado ``Plataforma para la Gestión Integral y Organización de Eventos mediante
Marketplace, Sistemas de Recomendación y Gemelos Digitales'', se
utilizaron herramientas de Inteligencia Artificial generativa como apoyo
en distintas etapas del desarrollo de la investigación.

Principalmente, se utilizó Claude, desarrollado por Anthropic, para
colaborar en la organización y estructuración de los capítulos, mejorar
la redacción de determinados apartados, revisar la claridad de los textos
y facilitar la comprensión de algunos conceptos técnicos relacionados con
las tecnologías utilizadas en el proyecto.

También se empleó como herramienta de apoyo para la elaboración de
encuestas y entrevistas, la organización de información, el análisis de
alternativas tecnológicas y la preparación de cálculos y estimaciones
correspondientes a la evaluación económica del proyecto.

El uso de estas herramientas tuvo como finalidad complementar el trabajo
realizado, facilitar determinadas tareas y mejorar la presentación de la
información. Las propuestas y respuestas generadas mediante Inteligencia
Artificial fueron utilizadas como material de apoyo, sujetas a revisión y
adaptación según las necesidades del proyecto.

Las decisiones relacionadas con el planteamiento del problema, los
objetivos, la metodología, la selección de tecnologías y el desarrollo de
la solución propuesta corresponden al autor del trabajo. Asimismo, la
responsabilidad sobre el contenido presentado, la validez de los datos
utilizados, los resultados y las conclusiones de la investigación recae
exclusivamente en el autor.

De esta manera, la Inteligencia Artificial fue utilizada como una
herramienta complementaria durante el proceso de elaboración de la tesis,
sin reemplazar el análisis, el criterio ni la responsabilidad académica
del estudiante.
